\section{Related Work}
\label{sec:related}

\para{LLM post-processing for diarization.} The closest prior work is \citet{wang2024diarizationlm}, who propose DiarizationLM for speaker diarization post-processing. Their method represents ASR and diarization outputs as speaker-tagged text, uses an LLM to correct speaker assignments, and applies transcript-preserving speaker transfer so the final transcript keeps the original words. Their results show large WDER reductions after fine-tuning, but also warn that zero-shot or lightly prompted models can delete, hallucinate, or fail to improve transcripts. We follow the transcript-preserving principle, but deliberately test the lightweight case: no fine-tuning, no word edits, and API-based speaker-ID-only outputs.

\para{Language-assisted diarization.} \citet{park2023enhancing} combine acoustic diarization evidence with lexical probabilities from language models inside contextual beam search, reporting strong relative improvements on speaker-attributed WER. \citet{india2025language} similarly show that linguistic cues can help speaker diarization in structured telephone interviews. These approaches support our hypothesis that text contains speaker-attribution information. They also motivate our negative result: lexical context alone is weaker when the candidate speakers are anonymized and the window is short.

\para{Speech-LLM systems for multi-speaker ASR.} Several recent systems move beyond text-only correction by conditioning LLMs on audio, diarization, speaker embeddings, or time segments. Serialized-output training with LLM decoders has improved multi-talker ASR on LibriMix and \amisdm~\cite{shi2024advancing}. Diarization-aware multi-speaker ASR feeds speaker/time triplets to an LLM backend~\cite{lin2025diarizationaware}, while DM-ASR formulates diarization-guided transcription as multi-turn dialogue generation~\cite{li2026dmasr}. SpeakerLM and JEDIS-LLM further explore end-to-end or streamable joint ASR and diarization with speaker registration and prompt-cache mechanisms~\cite{yin2025speakerlm,shi2025train}. Unlike these systems, our pilot does not train a speech model or expose acoustic embeddings to the LLM. We isolate what a prompted text LLM can do as a low-cost correction layer.

\para{Evaluation and toolkits.} Meeting transcription commonly reports word-level speaker-attribution metrics such as WDER, cpWER, and tcpWER. MeetEval standardizes these metrics for meeting transcription systems~\cite{vonneumann2023meeteval}. pyannote.audio provides widely used neural diarization building blocks and speaker diarization pipelines~\cite{bredin2020pyannote,bredin2023pyannote}. Our experiment uses local diarization annotations as ground truth and computes a small-speaker exact-permutation cpWER-style score locally because the resource-gathering phase could not install MeetEval under the available compiler constraints. This makes our cpWER values useful for within-pilot comparisons, not a substitute for a full MeetEval benchmark.

\para{Datasets.} \amisdm is a standard meeting benchmark derived from the AMI Meeting Corpus~\cite{carletta2006ami}. It is important for this study because it has four-speaker meeting-style interactions and high overlap rates in the selected shard. \simsamu provides a smaller two-speaker clinical-dialogue sanity dataset~\cite{medkit2023simsamu}. Using both lets us compare the easier two-speaker setting against the many-speaker condition that motivated the study.
